% Requires in preamble: \usepackage{booktabs, amsmath, cleveref}
\section{Residential Heating \& Cooling}\label{sec:heating_cooling}
Belgium's residential stock remains heavily fossil-based for space heating, but a
gradual electrification is under way through heat pumps, photovoltaics and, to a
lesser extent, active cooling. These technologies are key boundary conditions for
residential demand modelling, as they reshape both the annual energy balance and
the timing of electricity demand. This section characterises the carriers,
emitters and appliances that the technology layer of the building-stock model
assigns to the archetypes of \Cref{sec:size_compo}. It should be noted, however,
that recent Belgian heating-technology data are difficult to obtain: they are
often only available at regional level, are partly non-public, and are not always
nationally consistent \cite{rap_belgium_2024}.

\subsection{Heating Technology}
Space heating in Belgium is dominated by natural gas, followed by heating oil,
with electricity and other carriers (mainly biomass and district heat) making up
the remainder. The regional distribution of the main heating carrier is
summarised in \Cref{tab:heating_carrier_shares}, based on the Belgian Social
Climate Plan preparation report (EU-SILC 2022) \cite{socialclimateplan_be_2023}.
The regional contrast is pronounced: gas reaches 86\% of dwellings in the dense,
urban Brussels-Capital Region and 71\% in Flanders, but only 42\% in the more
rural Walloon Region, where heating oil remains prevalent (38\%) against 15\% in
Flanders and under 8\% in Brussels. A national cross-check confirms the
pattern, with fossil gas used by about 66\% of households for space heating and
59\% for domestic hot water, and heating oil the second most common space-heating
carrier \cite{rap_belgium_2024}.

\begin{table}[ht]
  \centering
  \caption{Regional distribution of the main residential heating carrier (\% of
           dwellings), EU-SILC 2022 \cite{socialclimateplan_be_2023}. The Walloon
           row sums to 99.7\% as published.}
  \label{tab:heating_carrier_shares}
  \begin{tabular}{lrrrrr}
    \toprule
    Region & Electricity & Natural gas & Heating oil & Other & Sum \\
    \midrule
    Flanders & 7.5  & 71.2 & 14.7 & 6.6  & 100.0 \\
    Brussels & 6.0  & 86.0 & 7.5  & 0.5  & 100.0 \\
    Wallonia & 6.0  & 42.4 & 37.5 & 13.8 & 99.7  \\
    \midrule
    Belgium  & 6.9  & 63.7 & 21.1 & 8.3  & 100.0 \\
    \bottomrule
  \end{tabular}
\end{table}

On the emission side, the existing stock relies overwhelmingly on high-temperature
hydronic radiators fed by gas or oil boilers. This is a central constraint for
electrification, since heat pumps operate efficiently only at low supply
temperatures. In the model, the as-is stock is therefore represented by these
lumped boiler-plus-radiator systems, whereas every renovated dwelling is assigned
a low-temperature package with 45\,\textdegree C radiators, consistent with the
EPBD classification of low-temperature emitters \cite{epbd_2024}.

\subsection{Heat Pumps}
Electrified heating is gaining ground. Heat pumps are now the second most common
technology in new heating-appliance sales, behind gas boilers, and the most
common type for space heating is air-to-water \cite{rap_belgium_2024}. The Joint
Research Centre reports that, in 2022, Belgium had an installed stock of roughly
209{,}000 heat pumps against about 2.8 million households using gas boilers, with
some 60{,}000 units sold that year ($+$118\% year-on-year); heat pumps
nonetheless represented only 11.8\% of combined heat-pump and boiler sales
\cite{JRC2022HeatPumpBelgium}. This electrification is reinforced by a rapidly
growing rooftop-solar base: cumulative PV capacity reached 11.7\,GW at the end of
2024, generating 8.32\,TWh and covering at least 10\% of national electricity
consumption, still dominated by sub-10\,kW residential systems
\cite{vermang2024_pvps_belgium}. Reflecting these trends, the model assigns every
renovated dwelling a package built around an air-to-water heat pump with
heat-pump domestic hot water and balanced heat-recovery ventilation, while as-is
dwellings retain their regional fossil-carrier mix.

\subsection{Cooling \& Reversible Heat Pumps}
Active cooling is still limited in the Belgian residential stock. In 2020, 88\% of
households had neither fixed nor mobile air-conditioning
\cite{fps_economy_household_energy_2020}, although ownership is rising: integrated
air-conditioning grew from 3\% of households in 2018 to 10\% in 2024, with
regional 2024 shares of 13\% in Flanders, 7\% in Wallonia and 4\% in Brussels
\cite{statbel_hbs_2024}. A structural driver of future cooling demand is the
diffusion of reversible heat pumps: the air-to-water units assigned to renovated
dwellings are technically capable of providing cooling, even though this
capability does not imply that cooling is actually operated. The model captures
this distinction explicitly, separating a heat pump's reversible-cooling
capability from active-cooling adoption, and explores the latter through scenarios
anchored on the current regional ownership shares. Under a warming climate, this
latent cooling capacity is a key uncertainty for both peak electricity demand and
the annual energy balance.
